\section{Conceptual framework} \label{sec:cf}

Suppose the firm has a collection of tasks that want performed which they can order for most to least valuable.
Let $v(L)$ be the marginal value of the task completed when $L$ efficiency units of labor are applied to the project.
If the firm hires a worker with productivity $y$ who works $h$ hours, $L = hy$.
The firm's decision of how many hours to buy is
\begin{align}
  \arg\max_{h^*} \quad \int_0^{h^*} v(hy) dh - wh^{*}  
\end{align} 
where $w$ is the worker's wage.
The firm's first order condition is simply $v(h^*y) = w$, i.e., where the value of output from the marginal hour is the wage.
This implies that $\epsilon^h_y = -1$, or that hiring a more productive worker simply means that fewer hours are purchased, even at the same price.
Suppose that worker offers a value multiplier on $e^{q_i}$.
The first order condition is now
\begin{align}
  e^{q} v(h y) = w.
\end{align} 
We have
\begin{align}
  h'(q) = - \frac{v(hy)}{y v'(hy)}
\end{align} 
\begin{align}
  \frac{h'(q)}{h(q)} = - \frac{v(hy)}{hy v'(hy)} > 0. 
\end{align}
\begin{align}
  q \frac{h'(q)}{h(q)} = - q \frac{1}{\epsilon^v_L}.
\end{align}



Suppose a firm has a project of size $Y$, which it can sell in the product market for $pY$.
If you hire a more productive worker, hours-worked goes down.
If you hire a worker with more surplus, the
Suppose the firm hires a worker with productivity technical productivity $y$, but who has latent quality effect $e^q$.
The firm's payoff is $pYe^{q_i} - w_i \frac{Y}{y}$.


The equilibrium in matching markets is profoundly shaped by search costs \citep{rogerson2005search}.
Much of these search costs are incurred by learning about the attributes of the other side of the market.
For example, job-seekers must learn about the existence of job openings, their individual probability of being hired, and at what wage.
Firms have to learn about the productivity and preferences of individual applicants. 
Some of the information needed in this match-making is provided willingly and credibly by market participants.
In particular, ``horizontal'' information such as the kind of work that is to be done, the location, the name of the firm and so on is typically provided truthfully.
Although this information is cheap talk, it works to coordinate the two sides and firms and workers have no strategic incentive to mislead \citep{farrell1996cheap}.
Other kinds of information such as educational background, criminal record or employment history are generally reported truthfully, even though a believed lie could be helpful, as third parties stand ready to verify that information \citep{autor2008economics}. 

Other kinds of information are important, but would not be viewed as credible if simply stated, and so they are signaled \citep{spence1973job}.  
In the classic case of education in the labor market, only high ability workers would choose to get costly education, and so a firm views education as a signal of ability.
In this case, the supply side of the market is vertically differentiated and the high-types want a separating equilibrium in which they are paid for their true productivity. 
Although it is the case that all else equal, the firm would prefer a more productive worker over a less productive worker at the same price, as in product markets, firms have vertical preferences.

Several recent papers have considered the effect of signaling in matching markets, in settings different to ours.
For example, \cite{lee2011propose} and \cite{coles2013preference} both consider the introduction of market interventions that make it easier for one side of the market to signal the other (participants in dating markets in the former, applicants in the academic job market in the latter).
In this case the signaling is a matter of horizontal differentiation.
By contrast, our paper focuses on signaling preferences over a vertically differentiated dimension, which brings with it the attendant risks of price markups and adverse selection as noted above.  

Most closely related to our paper is \cite{tadelis2011information}, who consider information disclosure in the wholesale auction market for used cars.
However, the signal is sent by the platform about the attributes of a good;
our signal is sent by the buyer about their preferences. 
In their experimental setting, information about used cars is {\em exogenously} randomly disclosed to market participants, before they choose which use car auctions to join.
This information disclosure induces a sorting effect that increases competition among similar-minded bidders for each type of car, and increases expected revenues---most strongly for cars at the extremes of the quality spectrum.

There are two major differences in our setting from \cite{tadelis2011information}.
First, in many markets honest information disclosure cannot be compelled (not least because the market operator may not possess the means or information to ensure it).
We focus on a setting where employers must {\em endogenously} choose whether or not to honestly signal workers about their preferences, and we find honest disclosure as a strong empirical result.
Second, the results of \cite{tadelis2011information} suggest that we should uniformly see {\em lower} wage bids in the sorted marketplace; however, we see wage markups in high tier jobs. 
As noted above, this is because the relative balance of demand and supply in the marketplace shifts bargaining power in the high tier from employers to workers. 

\cite{tadelis2011information}

\cite{anand2011} explores the role of advertising in communicating product attribtues to consumers. 
The context is television advertising for other television programs.
The authors fit a structural model, and find that exposure to a single advertisement reduces the probability of that viewer viewing their best option. 

Competitive information disclosure.
\cite{board2015competitive}

\cite{board2009}

\cite{budish2011}
\cite{cho2014}
\cite{day1976}
\cite{Silva2008}
\cite{Diamond1991}
\cite{Genesove1993}
\cite{gilson1984}
\cite{grossman1981}
\cite{healy2001}
\cite{jin2003}
\cite{johnson2006}
\cite{levin1996}
\cite{lewis2011}
\cite{milgrom1981}
\cite{Milgrom1982}
\cite{overby2009}
\cite{porter1995}
\cite{roberts2014}

\cite{coles2009signaling}
\cite{lee2009interviewing}
\cite{ely2009hiring}
\cite{ely2013adverse}
\cite{lee2011propose}
\cite{kushnir2013harmful}
